\clearpage
\fancyhead[L]{\small OBSERVATIONS.}
\SourceLeaf{468}
\begin{center}
\rule{0.55\linewidth}{0.4pt}\\[2pt]
Astronomische Nachrichten. Vol. X. No. 225. p. 143. 3 July 1832.\\[2pt]
\rule{0.55\linewidth}{0.4pt}
\end{center}

\begin{center}[\textsc{Gauss} to \textsc{Schumacher.}]\quad Göttingen, 12 May 1832.\end{center}

{\rightskip=0pt plus 1em At the last transit of Mercury, I missed the first contact at ingress because of a cloud and the second by chance. I did observe the egress\par}

\begin{center}
\begin{tabular}{ll}
inner contact at egress \mbox{.\,.\,.\,.\,.\,.\,.} & $4^{\rm u}25^{\rm m}32^{\rm s}1$ \\
outer contact \mbox{.\,.\,.\,.\,.\,.\,.\,.\,.\,.\,.\,.\,.\,.} & $4\;28\;22.5$ \\
\end{tabular}
\end{center}

I also observed its culmination with the meridian circle

\medskip
\noindent May 5. R.A.\quad $2^{\rm u}49^{\rm m}50^{\rm s}36$\par
\noindent $\delta+16^\circ28'43^{\rm s}6$ zero point determined from the nadir.\par
\noindent\hspace*{1em}---\quad ---\quad $44.4$ zero point from $\alpha$ Can. min. with the decl. of $\alpha$ from \textsc{Encke}’s \mbox{I. B.}

